\chapter{Beobachter}

\section{Tabellen}

\subsection{Einfache Tabelle}

\begin{center}
\begin{tabular}{|c|c|c|}
\hline
\textbf{Name} & \textbf{Fach} & \textbf{Note} \\
\hline
Alice & Mathematik & 1,0 \\
Bob & Physik & 1,5 \\
Charlie & Informatik & 1,3 \\
\hline
\end{tabular}
\end{center}

\subsection{Größere Tabelle mit Beschreibung}

\begin{center}
\begin{tabular}{|l|r|r|r|}
\hline
\textbf{Monat} & \textbf{Einnahmen} & \textbf{Ausgaben} & \textbf{Gewinn} \\
\hline
Januar & 10.000€ & 6.000€ & 4.000€ \\
Februar & 12.000€ & 7.500€ & 4.500€ \\
März & 15.000€ & 8.000€ & 7.000€ \\
April & 18.000€ & 9.500€ & 8.500€ \\
\hline
\end{tabular}
\end{center}

\section{Aufzählungen}

\subsection{Ungeordnete Liste}

\begin{itemize}
    \item Erster Punkt der Liste
    \item Zweiter Punkt mit \textbf{Formatierung}
    \begin{itemize}
        \item Verschachtelter Punkt 1
        \item Verschachtelter Punkt 2
    \end{itemize}
    \item Dritter Punkt
\end{itemize}

\subsection{Geordnete Liste}

\begin{enumerate}
    \item Erste Schritt im Prozess
    \item Zweite Schritt
    \begin{enumerate}
        \item Unterschritt A
        \item Unterschritt B
    \end{enumerate}
    \item Dritte Schritt
    \item Vierte Schritt
\end{enumerate}
